\Askhsh{22396_SOLUTION}{1}
\begin{ylh}{1}
    \item [ΛΥΣΗ]
\end{ylh}
\begin{alist}
    \item 
    \begin{rlist}
        \item Το τρίγωνο $ΑΔΒ$ είναι ορθογώνιο με $\hat{\Delta} = 90^\circ$,
        $ΑΒ = ΑΓ = ΑΔ + ΔΓ = 3 + 2 = 5$
        και $ΑΔ = 3$.
        Επομένως από το Πυθαγόρειο θεώρημα θα έχουμε
        \[ ΒΔ^2 = ΑΒ^2 - ΑΔ^2 \]
        \[ ΒΔ^2 = 5^2 - 3^2 \]
        \[ ΒΔ^2 = 16 \]
        $ΒΔ = 4$.
        
        \item Το εμβαδόν του τριγώνου $ΑΒΓ$ είναι
        \[ (ΑΒΓ) = \frac{ΑΓ \cdot ΒΔ}{2} = \frac{5 \cdot 4}{2} = 10. \]
    \end{rlist}
    
    \item 
    \begin{rlist}
        \item Το τρίγωνο $ΑΒΕ$ είναι ορθογώνιο με $Α\hat{Β}Ε = 90^\circ$, άρα για το ύψος του $ΒΔ$ που αντιστοιχεί στην υποτείνουσα $ΑΕ$, θα έχουμε
        \[ ΒΔ^2 = ΔΑ \cdot ΔΕ \]
        \[ 4^2 = 3 \cdot ΔΕ \]
        \[ ΔΕ = \frac{16}{3}. \]
        
        \item Είναι $ΓΕ = ΔΕ - ΔΓ = \frac{16}{3} - 2 = \frac{10}{3}$. Το εμβαδόν του τριγώνου $ΒΓΕ$ είναι
        \[ (ΒΓΕ) = \frac{ΓΕ \cdot ΒΔ}{2} = \frac{\frac{10}{3} \cdot 4}{2} = \frac{20}{3}. \]
    \end{rlist}
\end{alist}